\def\level{Première}  % Classe à droite
\def\course{NSI}
\def\chaptername{La langage assembleur} % Titre au centre
\def\docname{ARSE01-ACT1}        % Identifiant à gauche

\pagestyle{cours_FP}
\makeheader{} % Affiche le bandeau

\begin{easybox}{Les registres}{}

    \begin{itemize}
        \item un \textbf{registre} est un emplacement de mémoire interne à un processeur. Les registres sont les mémoires les plus rapides d'un ordinateur mais dont le coût de fabrication est le plus élevé.
        \item les registres sont utilisés pour stocker des opérandes et des résultats intermédiaires lors des opérations effectuées dans l'UAL (unité arithmétique et logique) du processeur.
        \item la plupart des PC actuels ont des registres de taille 64 bits.
    \end{itemize}

        
\end{easybox}


\begin{easybox}{Initiation à l'assembleur}{}
	Le langage \textbf{assembleur} est un langage de programmation très particulier. Il possède des instructions:
	\begin{itemize}
		\item arithmétiques: addition, soustraction, multiplication \ldots
		\item transfert de données: déplace des données du registre vers la mémoire et inversement
		\item rupture de séquence: des sauts dans le programme pour aller à une ligne particulière.
	\end{itemize}
\end{easybox}

\begin{easybox}{Simulateur assembleur}{}
	 \href{https://www.peterhigginson.co.uk/AQA/}{https://www.peterhigginson.co.uk/AQA/}
\end{easybox}



\begin{exemple}[Voici un petit programme qui effectue une addition et stocke le résultat.]{}
	\begin{verbatim}
		MOV R0, #22            # Stocke la valeur 22 dans le registre R0
		MOV R1, #36            # Stocke la valeur 36 dans le regsitre R1
		ADD R2, R0, R1         # Addition R0 et R1 et mets le résultat dans R2
		STR R2, 24             # Stocke R2 à l'adresse mémoire 24
	\end{verbatim}
\end{exemple}

\renewcommand{\arraystretch}{1.21}
\begin{longtable}{|l|l|l|l|}
    \hline
    \textbf{Commande}  & \textbf{Description}  & \textbf{Exemple}  & \\    \hline
    \endfirsthead  \hline
    \endfoot  \hline
    charger   & \texttt{LDR R1, 78} & \texttt{R1 <- @78} & Place la valeur stockée à l'adresse 78 dans \texttt{R1}\\
    stocker   & \texttt{STR R3, 125} & \texttt{@125 <- R3} & Place la valeur du registre \texttt{R3} à l'adresse 125 \\
    déplacer  & \texttt{MOV R0, \#23} & \texttt{R0 <- 23} & Place la valeur 23 dans \texttt{R0} \\ \hline
    ajouter   & \texttt{ADD R1, R0, \#128} &\texttt{R1 <- R0 + 128}& Additionne 128 à \texttt{R0} et le place dans \texttt{R1} \\
    ajouter   & \texttt{ADD R0, R1, R2} & \texttt{R0 <- R1 + R2} & Additionne \texttt{R1} et \texttt{R2} et le place dans \texttt{R0} \\ 
    soustraire & \texttt{SUB R1, R0, \#128} & \texttt{R1 <- R0 - 128} & Soustrait 128 de \texttt{R0} et le place dans \texttt{R1} \\
    soustraire & \texttt{SUB R0, R1, R2} & \texttt{R0<- R1 - R2} & Soustrait \texttt{R2} de \texttt{R1} et le place dans \texttt{R0} \\ \hline
    arrêt & \texttt{HALT} & &Arrêt de l'exécution du programme \\
    saut  & \texttt{B label} &  & La prochaine instruction se situe en \texttt{label}\\
    comparer & \texttt{CMP R0, \#23} & \texttt{R0 == 23} & Compare \texttt{R0} au nombre 23 - suivi par ``saut''\\
    saut ``égal'' & \texttt{BEQ label} & \texttt{if R0 == 23}& Si le dernier \verb+CMP+ est égal, saute à \texttt{label}\\
    saut ``différent'' & \texttt{BNE label} & \texttt{if R0 != 23}& Si le dernier \texttt{CMP} est différent, saute à \texttt{label}\\
    saut ``plus grand'' & \texttt{BGT label} & \texttt{if R0 > 23}& Si le dernier \texttt{CMP} est plus grand, saute à \texttt{label}\\
    saut ``plus petit'' & \texttt{BLT label} & \texttt{if R0 < 23}& Si le dernier \texttt{CMP} est plus petit, saute à \texttt{label}\\
	\hline
	entrée clavier & \texttt{INP R0} & \texttt{R0 = input()} & Demande une entrée clavier et stocke dans \texttt{R0} \\
	sortie affichage & \texttt{OUT R0} & \texttt{print(R0)} & Affiche le contenu de \texttt{R0} \\
	\hline
	label & \texttt{mon\_label} & & crée un label dans le programme pour le boucles \\
	constante & \texttt{mavaleur: 7} & \texttt{mavaleur = 7} & crée une constante de valeur 7\\ \hline
\end{longtable}


\newpage


\begin{minipage}[t]{0.48\linewidth}
	\begin{exercice_cours}{Expliquer ce que fait cet algorithme:}{}
		
		\begin{verbatim}
			MOV R0, #10
			MOV R1, #20
			ADD R2, R1, R0
			ADD R2, R2, #5
		\end{verbatim}
	\end{exercice_cours}
\end{minipage}\hfill
\begin{minipage}[t]{0.48\linewidth}
	\begin{exercice_cours}{Expliquer ce que fait cet algorithme:}{}
		
		\begin{verbatim}
			INP R0
			ADD R0, R0, #1
			OUT R0
			HALT
		\end{verbatim}
	\end{exercice_cours}
\end{minipage}



\begin{minipage}[t]{0.48\linewidth}
	\begin{exercice_cours}{Expliquer ce que fait cet algorithme:}{}
		
		\begin{verbatim}
     INP R0
     INP R1
     MOV R2, #0
boucle:
     CMP R1, #0
     BEQ fin
     SUB R1, R1, #1
     ADD R2, R2, R0
     B boucle
fin:
     OUT R2
     HALT
\end{verbatim}
	\end{exercice_cours}

	\begin{exercice_cours}{Programmer en Assembleur}
		\begin{lstlisting}
x = 0
while x < 3:
	x = x + 1
print(x)
		\end{lstlisting}
	\end{exercice_cours}
\end{minipage}\hfill
\begin{minipage}[t]{0.48\linewidth}
	\begin{exercice_cours}{Expliquer ce que fait cet algorithme:}{}
		
		\begin{verbatim}
     MOV R12, #999
     MOV R11, #888
     INP R0 
     STR R0, 40 
boucle:
     CMP R0, #2 
     BLT decision
     SUB R0, R0, #2
     B boucle
decision:
     CMP R0, #1
     BEQ label
     OUT R12
     B halte
label:
     OUT R11
halte:
     HALT
	  
		\end{verbatim}
	\end{exercice_cours}
\end{minipage}

\begin{minipage}[t]{0.48\linewidth}
	\begin{exercice_cours}{}{}
		Écrire en assembleur un programme qui calcule la somme des $n$ premiers entiers naturels: $1+2+3+\ldots+n$
		
	\end{exercice_cours}
\end{minipage}\hfill
\begin{minipage}[t]{0.48\linewidth}
	\begin{exercice_cours}{}{}
Écrire un programme qui demande 2 entiers positifs et affiche le plus grand
	\end{exercice_cours}
\end{minipage}

\begin{minipage}[t]{0.48\linewidth}
	\begin{exercice_cours}{}{}
		Écrire en assembleur un programme qui demande un entier naturel $n$ et affiche $2^n$
		
	\end{exercice_cours}
\end{minipage}\hfill
\begin{minipage}[t]{0.48\linewidth}
	\begin{exercice_cours}{}{}
Écrire un programme qui demande trois entiers positifs et les affiche par ordre croissant	\end{exercice_cours}
\end{minipage}

%%

\begin{minipage}[t]{0.48\linewidth}
	\begin{exercice_cours}{Traduire en assembleur}{}
\begin{lstlisting}
capital = 2000
for annee in range(10):
	capital = capital + 150
print(capital)
\end{lstlisting}		
	\end{exercice_cours}
\end{minipage}\hfill
\begin{minipage}[t]{0.48\linewidth}
	\begin{exercice_cours}{Traduire en assembleur}{}

\begin{lstlisting}
capital = 2000
n = 0
while capital < 3000:
	capital = capital + 123
print(n)
			\end{lstlisting}
		\end{exercice_cours}
	\end{minipage}


